% =====================================================================
%  PROYECTO DISEÑO EN ACERO (CIV-336) - USM - 2do SEMESTRE 2026
%  GRUPO 5  -  TAREA 1: Cargas, Modelación, Análisis Plástico y Tracción
%  Documento autocontenido: compila directamente en Overleaf (pdfLaTeX).
%  Las figuras están hechas en TikZ/pgfplots (no requiere imágenes).
% =====================================================================
\documentclass[11pt,letterpaper]{article}
\usepackage[utf8]{inputenc}
\usepackage[T1]{fontenc}
\usepackage[spanish,es-noshorthands,es-tabla]{babel}
\usepackage{lmodern}
\usepackage[margin=2.5cm]{geometry}
\usepackage{amsmath,amssymb}
\usepackage{booktabs,array,multirow,tabularx}
\usepackage{float}
\usepackage{caption}
\usepackage{xcolor}
\usepackage{tikz}
\usetikzlibrary{arrows.meta,patterns,calc,decorations.pathmorphing}
\usepackage{pgfplots}
\pgfplotsset{compat=1.16}
\usepackage{fancyhdr}
\usepackage[hidelinks]{hyperref}
\setlength{\parskip}{0.5em}
\setlength{\headheight}{14pt}
\setlength{\parindent}{0pt}
\pagestyle{fancy}
\fancyhf{}
\lhead{CIV-336 Diseño en Acero}
\rhead{Proyecto -- Tarea 1 -- Grupo 5}
\cfoot{\thepage}
\renewcommand{\arraystretch}{1.15}
\newcommand{\tonf}{\,\mathrm{tonf}}
\newcommand{\tonfm}{\,\mathrm{tonf\cdot m}}

\begin{document}

% ------------------------------------------------------------------ PORTADA
\begin{titlepage}
\centering
{\large Universidad Técnica Federico Santa María\par}
{\large Departamento de Obras Civiles\par}
\vspace{3cm}
{\Large CIV-336 Diseño en Acero\par}
\vspace{1cm}
{\huge\bfseries Proyecto Semestral\par}
\vspace{0.4cm}
{\LARGE Tarea 1: Cargas, Modelación, Análisis Plástico\\ y Diseño a Tracción\par}
\vspace{2cm}
{\Large\bfseries Grupo 5\par}

\vspace{0.3cm}
{\large $L = 10\ \mathrm{m}$ \quad -- \quad $H = 3{,}0\ \mathrm{m}$\par}
\vfill
\begin{flushleft}
\begin{tabular}{ll}
\textbf{Integrantes:} & Nombre Apellido 1 -- Rol \\
                      & Nombre Apellido 2 -- Rol \\
                      & Nombre Apellido 3 -- Rol \\[0.3cm]
\textbf{Profesor:}    & Nombre del profesor \\
\textbf{Fecha:}       & \today
\end{tabular}
\end{flushleft}
\end{titlepage}

\tableofcontents
\newpage

\section{Introducción y descripción de la estructura}

La presente memoria de cálculo corresponde a la Tarea 1 del proyecto de la asignatura
CIV-336 Diseño en Acero. Se estudia un edificio de cinco pisos estructurado en base a marcos de acero,
con marcos arriostrados concéntricos en los ejes A, B, 1 y 5, y marcos rígidos (resistentes a momento)
en los ejes 2, 3 y 4. En esta entrega se presentan las bases de cálculo, la determinación de las cargas,
la descripción del modelo realizado en SAP2000, el resumen de esfuerzos internos, el análisis plástico del
eje 2 y el diseño a tracción de la riostra más traccionada.

Las dimensiones de la estructura dependen del número de grupo (Tabla 2 del enunciado). Para el
\textbf{Grupo 5} se tiene $L = 10\ \mathrm{m}$ y $H = 3{,}0\ \mathrm{m}$, de donde resulta la geometría
de la Tabla~\ref{tab:geom}.

\begin{table}[H]
\centering
\caption{Geometría de la estructura (Grupo 5).}
\label{tab:geom}
\begin{tabular}{lcc}
\toprule
Descripción & Expresión & Valor \\
\midrule
Vanos extremos (ejes 1--2 y 4--5) & $0{,}75L$ & $7{,}50$ m \\
Vanos centrales (ejes 2--3 y 3--4) & $L$ & $10{,}00$ m \\
Largo total en dirección X & $3{,}5L$ & $35{,}00$ m \\
Ancho en dirección Y (ejes A--B) & $0{,}8L$ & $8{,}00$ m \\
Altura primer piso & $1{,}1H$ & $3{,}30$ m \\
Altura pisos 2 a 5 & $H$ & $3{,}00$ m \\
Altura total $H_{total}$ & $1{,}1H + 4H$ & $15{,}30$ m \\
Área de planta por piso & $3{,}5L \times 0{,}8L$ & $280{,}0\ \mathrm{m^2}$ \\
\bottomrule
\end{tabular}
\end{table}

\begin{figure}[H]
\centering
\begin{tikzpicture}[x=0.4cm,y=0.4cm,font=\small]
\draw[gray,dashed] (0,-1.5) -- (0,9.5);
\draw[gray,dashed] (7.5,-1.5) -- (7.5,9.5);
\draw[gray,dashed] (17.5,-1.5) -- (17.5,9.5);
\draw[gray,dashed] (27.5,-1.5) -- (27.5,9.5);
\draw[gray,dashed] (35,-1.5) -- (35,9.5);
\draw[gray,dashed] (-1.5,0) -- (36.5,0);
\draw[gray,dashed] (-1.5,8) -- (36.5,8);
\node[draw,circle,inner sep=1.5pt] at (0,10.6) {1};
\node[draw,circle,inner sep=1.5pt] at (7.5,10.6) {2};
\node[draw,circle,inner sep=1.5pt] at (17.5,10.6) {3};
\node[draw,circle,inner sep=1.5pt] at (27.5,10.6) {4};
\node[draw,circle,inner sep=1.5pt] at (35,10.6) {5};
\node[draw,circle,inner sep=1.5pt] at (-2.6,0) {A};
\node[draw,circle,inner sep=1.5pt] at (-2.6,8) {B};
\draw[thick] (0,0) rectangle (35,8);
\draw[thick] (0,0) -- (0,8);
\fill (0,0) rectangle +(0.35,0.35); \fill (0,0) rectangle +(-0.35,-0.35); \fill (0,0) rectangle +(0.35,-0.35); \fill (0,0) rectangle +(-0.35,0.35);
\fill (0,8) rectangle +(0.35,0.35); \fill (0,8) rectangle +(-0.35,-0.35); \fill (0,8) rectangle +(0.35,-0.35); \fill (0,8) rectangle +(-0.35,0.35);
\draw[thick] (7.5,0) -- (7.5,8);
\fill (7.5,0) rectangle +(0.35,0.35); \fill (7.5,0) rectangle +(-0.35,-0.35); \fill (7.5,0) rectangle +(0.35,-0.35); \fill (7.5,0) rectangle +(-0.35,0.35);
\fill (7.5,8) rectangle +(0.35,0.35); \fill (7.5,8) rectangle +(-0.35,-0.35); \fill (7.5,8) rectangle +(0.35,-0.35); \fill (7.5,8) rectangle +(-0.35,0.35);
\draw[thick] (17.5,0) -- (17.5,8);
\fill (17.5,0) rectangle +(0.35,0.35); \fill (17.5,0) rectangle +(-0.35,-0.35); \fill (17.5,0) rectangle +(0.35,-0.35); \fill (17.5,0) rectangle +(-0.35,0.35);
\fill (17.5,8) rectangle +(0.35,0.35); \fill (17.5,8) rectangle +(-0.35,-0.35); \fill (17.5,8) rectangle +(0.35,-0.35); \fill (17.5,8) rectangle +(-0.35,0.35);
\draw[thick] (27.5,0) -- (27.5,8);
\fill (27.5,0) rectangle +(0.35,0.35); \fill (27.5,0) rectangle +(-0.35,-0.35); \fill (27.5,0) rectangle +(0.35,-0.35); \fill (27.5,0) rectangle +(-0.35,0.35);
\fill (27.5,8) rectangle +(0.35,0.35); \fill (27.5,8) rectangle +(-0.35,-0.35); \fill (27.5,8) rectangle +(0.35,-0.35); \fill (27.5,8) rectangle +(-0.35,0.35);
\draw[thick] (35,0) -- (35,8);
\fill (35,0) rectangle +(0.35,0.35); \fill (35,0) rectangle +(-0.35,-0.35); \fill (35,0) rectangle +(0.35,-0.35); \fill (35,0) rectangle +(-0.35,0.35);
\fill (35,8) rectangle +(0.35,0.35); \fill (35,8) rectangle +(-0.35,-0.35); \fill (35,8) rectangle +(0.35,-0.35); \fill (35,8) rectangle +(-0.35,0.35);
\draw[red,line width=1.6pt] (0,0) -- (7.5,0) (27.5,0) -- (35,0) (0,8) -- (7.5,8) (27.5,8) -- (35,8);
\draw[red,line width=1.6pt] (0,0) -- (0,8) (35,0) -- (35,8);
\draw[<->] (0,-1.0) -- node[below]{7,5 m} (7.5,-1.0);
\draw[<->] (7.5,-1.0) -- node[below]{10 m} (17.5,-1.0);
\draw[<->] (17.5,-1.0) -- node[below]{10 m} (27.5,-1.0);
\draw[<->] (27.5,-1.0) -- node[below]{7,5 m} (35,-1.0);
\draw[<->] (36.2,0) -- node[right]{8 m} (36.2,8);
\fill[blue] (17.5,4) circle (0.3) node[above right,blue]{CM};
\draw[->,thick] (-1.5,-3.5) -- (1.5,-3.5) node[right]{X};
\draw[->,thick] (-1.5,-3.5) -- (-1.5,-0.8) node[above]{Y};
\end{tikzpicture}
\caption{Vista en planta (Grupo 5). En rojo se indican los vanos arriostrados.}
\label{fig:planta}
\end{figure}

\begin{figure}[H]
\centering
\begin{tikzpicture}[x=0.205cm,y=0.205cm,font=\scriptsize]
\draw[thick] (0,0) -- (0,15.3);
\draw (0,0) -- ++(-0.7,-0.6) -- ++(1.4,0) -- cycle;
\draw[thick] (7.5,0) -- (7.5,15.3);
\draw (7.5,0) -- ++(-0.7,-0.6) -- ++(1.4,0) -- cycle;
\draw[thick] (17.5,0) -- (17.5,15.3);
\draw (17.5,0) -- ++(-0.7,-0.6) -- ++(1.4,0) -- cycle;
\draw[thick] (27.5,0) -- (27.5,15.3);
\draw (27.5,0) -- ++(-0.7,-0.6) -- ++(1.4,0) -- cycle;
\draw[thick] (35,0) -- (35,15.3);
\draw (35,0) -- ++(-0.7,-0.6) -- ++(1.4,0) -- cycle;
\draw[thick] (0,3.3) -- (35,3.3);
\draw[thick] (0,6.3) -- (35,6.3);
\draw[thick] (0,9.3) -- (35,9.3);
\draw[thick] (0,12.3) -- (35,12.3);
\draw[thick] (0,15.3) -- (35,15.3);
\draw[red,thick] (0,0) -- (7.5,3.3);
\draw[red,thick] (27.5,3.3) -- (35,0);
\draw[red,thick] (0,3.3) -- (7.5,6.3);
\draw[red,thick] (27.5,6.3) -- (35,3.3);
\draw[red,thick] (0,6.3) -- (7.5,9.3);
\draw[red,thick] (27.5,9.3) -- (35,6.3);
\draw[red,thick] (0,9.3) -- (7.5,12.3);
\draw[red,thick] (27.5,12.3) -- (35,9.3);
\draw[red,thick] (0,12.3) -- (7.5,15.3);
\draw[red,thick] (27.5,15.3) -- (35,12.3);
\node at (0,-2) {1};
\node at (7.5,-2) {2};
\node at (17.5,-2) {3};
\node at (27.5,-2) {4};
\node at (35,-2) {5};
\node at (17.5,-4) {\small Elevación ejes A y B};
\node[left] at (-0.5,3.3) {N1};
\node[left] at (-0.5,6.3) {N2};
\node[left] at (-0.5,9.3) {N3};
\node[left] at (-0.5,12.3) {N4};
\node[left] at (-0.5,15.3) {N5};
\draw[thick] (41,0) -- (41,15.3);
\draw (41,0) -- ++(-0.7,-0.6) -- ++(1.4,0) -- cycle;
\draw[thick] (49,0) -- (49,15.3);
\draw (49,0) -- ++(-0.7,-0.6) -- ++(1.4,0) -- cycle;
\draw[thick] (41,3.3) -- (49,3.3);
\draw[thick] (41,6.3) -- (49,6.3);
\draw[thick] (41,9.3) -- (49,9.3);
\draw[thick] (41,12.3) -- (49,12.3);
\draw[thick] (41,15.3) -- (49,15.3);
\draw[red,thick] (41,3.3) -- (49,0);
\draw[red,thick] (41,3.3) -- (49,6.3);
\draw[red,thick] (41,9.3) -- (49,6.3);
\draw[red,thick] (41,9.3) -- (49,12.3);
\draw[red,thick] (41,15.3) -- (49,12.3);
\node at (41,-2) {A}; \node at (49,-2) {B};
\node at (45,-4) {\small Ejes 1 y 5};
\draw[thick] (55,0) -- (55,15.3);
\draw (55,0) -- ++(-0.7,-0.6) -- ++(1.4,0) -- cycle;
\draw[thick] (63,0) -- (63,15.3);
\draw (63,0) -- ++(-0.7,-0.6) -- ++(1.4,0) -- cycle;
\draw[thick] (55,3.3) -- (63,3.3);
\draw[thick] (55,6.3) -- (63,6.3);
\draw[thick] (55,9.3) -- (63,9.3);
\draw[thick] (55,12.3) -- (63,12.3);
\draw[thick] (55,15.3) -- (63,15.3);
\node at (55,-2) {A}; \node at (63,-2) {B};
\node at (59,-4) {\small Ejes 2, 3 y 4};
\draw[<->] (64.2,0) -- node[right]{3,3 m} (64.2,3.3);
\draw[<->] (64.2,3.3) -- node[right]{3 m} (64.2,6.3);
\draw[<->] (64.2,6.3) -- node[right]{3 m} (64.2,9.3);
\draw[<->] (64.2,9.3) -- node[right]{3 m} (64.2,12.3);
\draw[<->] (64.2,12.3) -- node[right]{3 m} (64.2,15.3);
\end{tikzpicture}
\caption{Elevaciones de la estructura. En rojo, las riostras (diagonales).}
\label{fig:elev}
\end{figure}

\section{Bases de cálculo}

\subsection{Normas y documentos de referencia}
\begin{itemize}
  \item Enunciado del Proyecto -- Diseño en Acero (CIV-336), Segundo Semestre 2026.
  \item NCh1537.Of2009: Diseño estructural -- Cargas permanentes y cargas de uso.
  \item NCh3171.Of2010: Diseño estructural -- Disposiciones generales y combinaciones de carga.
  \item NCh433.Of1996 Mod.2009 (referencial, para el peso sísmico y la distribución de fuerzas).
  \item ANSI/AISC 360-16: \emph{Specification for Structural Steel Buildings} (método LRFD).
  \item AISC \emph{Steel Construction Manual}, 15\textsuperscript{a} edición (propiedades de perfiles).
\end{itemize}

\subsection{Materiales}
\begin{table}[H]
\centering
\caption{Materiales considerados.}
\begin{tabular}{llcc}
\toprule
Elemento & Material & $F_y$ [MPa] & $F_u$ [MPa] \\
\midrule
Vigas y columnas (perfiles W) & ASTM A572 Gr.50 & 345 & 450 \\
Riostras (HSS cuadrados) & ASTM A500 Gr.C & 345 & 427 \\
Planchas de conexión (gusset) & ASTM A36 & 248 & 400 \\
Soldadura & Electrodo E70XX & -- & $F_{EXX} = 482$ \\
Losa & Hormigón armado & \multicolumn{2}{c}{$\gamma = 2500\ \mathrm{kgf/m^3}$} \\
\bottomrule
\end{tabular}
\end{table}
Módulo de elasticidad del acero $E = 200\,000$ MPa, módulo de Poisson $\nu = 0{,}3$ y
peso unitario $\gamma_a = 7850\ \mathrm{kgf/m^3}$.

\subsection{Método de diseño}
Se utiliza el método de Diseño por Factores de Carga y Resistencia (LRFD) según AISC 360-16, con las
combinaciones de carga de la NCh3171 (sección~\ref{sec:combos}).

\section{Cargas}

\subsection{Cargas permanentes (D)}
\begin{table}[H]
\centering
\caption{Cargas permanentes (Tabla 1 del enunciado).}
\begin{tabular}{lcc}
\toprule
Descripción & Cálculo & Valor \\
\midrule
Peso propio acero estructural & $\gamma_a = 7850\ \mathrm{kgf/m^3}$ & SAP2000 (automático) \\
Conexiones & $+10\%$ del peso propio & factor $1{,}10$ \\
Losa de H.A. ($e = 18$ cm) & $0{,}18 \times 2500$ & $450\ \mathrm{kgf/m^2}$ \\
Tabiques, terminaciones, etc. & -- & $100\ \mathrm{kgf/m^2}$ \\
\midrule
Sobrecarga permanente (SCP) & losa + tabiques & $550\ \mathrm{kgf/m^2}$ \\
\bottomrule
\end{tabular}
\end{table}

El peso propio de los elementos de acero se considera en el patrón \texttt{PP} con un multiplicador de
peso propio igual a $1{,}10$, lo que incluye el $10\%$ adicional por conexiones. La losa y los tabiques
se consideran en el patrón \texttt{SCP}. El peso de acero resultante del prediseño es:
columnas $22{,}3\tonf$, vigas X $28{,}2\tonf$,
vigas Y $20{,}1\tonf$, riostras X $12{,}7\tonf$ y
riostras Y $8{,}5\tonf$ (incluye el $10\%$ de conexiones), con un total de
$91{,}7\tonf$.

\subsection{\texorpdfstring{Cargas de uso (L y L\textsubscript{r})}{Cargas de uso (L y Lr)}}
\begin{itemize}
  \item Pisos 1 a 4: sobrecarga de uso $q_L = 5\ \mathrm{kPa} = 510\ \mathrm{kgf/m^2}$
        (patrón \texttt{L}).
  \item Techo (nivel 5): techo transitable según NCh1537,
        $q_{L_r} = 2{,}0\ \mathrm{kPa} = 204\ \mathrm{kgf/m^2}$
        (patrón \texttt{LR}). % <<< VERIFICAR con la Tabla 4 de NCh1537 vista en clases
\end{itemize}
No se aplica reducción de la sobrecarga por área tributaria (criterio conservador).

\subsection{Distribución de las cargas de losa a las vigas}
La losa no se modela como elemento estructural; su peso, los tabiques y las sobrecargas se traspasan a las
vigas mediante áreas tributarias, considerando losas apoyadas en sus cuatro bordes (trabajo en dos
direcciones) con líneas de rotura a $45^\circ$. Así, en cada paño de $a \times b$ ($a \le b$):
las vigas del lado corto reciben una carga triangular de altura $q\,a/2$ y las del lado largo
una carga trapezoidal de altura $q\,a/2$. Para el Grupo 5 los paños son de
$7{,}5 \times 8{,}0$ m (vanos extremos) y $10{,}0 \times 8{,}0$ m (vanos centrales).

\begin{table}[H]
\centering
\caption{Cargas máximas (ordenada máxima) transmitidas por la losa a las vigas, por paño.}
\begin{tabular}{llccc}
\toprule
Paño & Viga & Forma & SCP [tonf/m] & L [tonf/m] \\
\midrule
$7{,}5 \times 8$ & Vigas X (luz 7,5 m) & triangular & $2{,}06$ & $1{,}91$ \\
$7{,}5 \times 8$ & Vigas Y (luz 8 m) & trapezoidal & $2{,}06$ & $1{,}91$ \\
$10 \times 8$ & Vigas X (luz 10 m) & trapezoidal & $2{,}20$ & $2{,}04$ \\
$10 \times 8$ & Vigas Y (luz 8 m) & triangular & $2{,}20$ & $2{,}04$ \\
\bottomrule
\end{tabular}
\end{table}
Las vigas Y de los ejes 2, 3 y 4 reciben la suma de los dos paños adyacentes. En el techo la carga
de uso es la indicada para $L_r$.

\subsection{Cargas sísmicas (E)}
De acuerdo al enunciado, el esfuerzo de corte basal se obtiene como $Q = 0{,}18\,P$, donde $P$ es el peso
sísmico de la estructura, igual a las cargas permanentes más un $25\%$ de la sobrecarga de uso
(también se incluye el $25\%$ de la sobrecarga del techo transitable). El peso de cada nivel considera la
SCP de la planta, el peso de las vigas del nivel y la mitad de las columnas y riostras de los pisos
superior e inferior. Las fuerzas se distribuyen en altura según:
\begin{equation}
F_k = \frac{A_k\,P_k}{\sum_{j=1}^{5} A_j\,P_j}\,Q,
\qquad
A_k = \sqrt{1-\frac{Z_{k-1}}{H_{total}}} - \sqrt{1-\frac{Z_k}{H_{total}}}
\end{equation}
con $H_{total} = 15{,}30$ m.

\begin{table}[H]
\centering
\caption{Peso sísmico y fuerzas sísmicas por nivel.}
\label{tab:sismo}
\begin{tabular}{cccccccc}
\toprule
Nivel & $Z_k$ [m] & Acero [tonf] & SCP [tonf] & $0{,}25\,L$ [tonf] & $P_k$ [tonf] & $A_k$ & $F_k$ [tonf] \\
\midrule
1 & $3{,}30$ & $18{,}5$ & $154{,}0$ & $35{,}7$ & $208{,}2$ & $0{,}1144$ & $22{,}11$ \\
2 & $6{,}30$ & $18{,}2$ & $154{,}0$ & $35{,}7$ & $207{,}9$ & $0{,}1186$ & $22{,}91$ \\
3 & $9{,}30$ & $18{,}2$ & $154{,}0$ & $35{,}7$ & $207{,}9$ & $0{,}1407$ & $27{,}18$ \\
4 & $12{,}30$ & $18{,}2$ & $154{,}0$ & $35{,}7$ & $207{,}9$ & $0{,}1834$ & $35{,}42$ \\
5 & $15{,}30$ & $14{,}0$ & $154{,}0$ & $14{,}3$ & $182{,}2$ & $0{,}4428$ & $74{,}94$ \\
\midrule
$\Sigma$ & & & & & $1014{,}2$ & $1{,}0000$ & $182{,}56$ \\
\bottomrule
\end{tabular}
\end{table}

\[
P = 1014{,}2\tonf \quad\Rightarrow\quad Q = 0{,}18 \times 1014{,}2 = 182{,}6\tonf
\qquad \left(\textstyle\sum A_j P_j = 196{,}58\tonf\right)
\]

Las fuerzas $F_k$ se aplican en el centro de masas (CM) de cada piso, ubicado en
$(X, Y) = (17{,}50;\ 4{,}00)$ m dada la simetría de la planta, en forma
independiente en las direcciones X (patrón \texttt{EX}) e Y (patrón \texttt{EY}).

\subsection{Combinaciones de carga}\label{sec:combos}
Se utilizan las combinaciones de la NCh3171 para el método LRFD, con $D = PP + SCP$. Al no existir
cargas de viento, nieve ni lluvia, éstas se omiten. Las cargas sísmicas se consideran con ambos signos
y en forma independiente en cada dirección. Se usa factor $1{,}0$ para la sobrecarga $L$ en las
combinaciones sísmicas (conservador).
\begin{table}[H]
\centering
\caption{Combinaciones de carga (NCh3171).}
\begin{tabular}{cl}
\toprule
Combinación & Expresión \\
\midrule
C1 & $1.4D$ \\
C2 & $1.2D + 1.6L + 0.5L_r$ \\
C3 & $1.2D + 1.6L_r + L$ \\
C4 & $1.2D + L + 1.4E_x$ \\
C5 & $1.2D + L - 1.4E_x$ \\
C6 & $1.2D + L + 1.4E_y$ \\
C7 & $1.2D + L - 1.4E_y$ \\
C8 & $0.9D + 1.4E_x$ \\
C9 & $0.9D - 1.4E_x$ \\
C10 & $0.9D + 1.4E_y$ \\
C11 & $0.9D - 1.4E_y$ \\
\midrule
ENVOLVENTE & Envolvente de C1 a C11 \\
\bottomrule
\end{tabular}
\end{table}

\section{Modelo estructural en SAP2000}

\subsection{Descripción del modelo}
Se realizó un modelo tridimensional en SAP2000 v27 (archivo \texttt{G5\_Edificio3D.sdb}, entregado también
en formato de texto \texttt{G5\_Edificio3D.s2k}). Las principales características y supuestos son:
\begin{itemize}
  \item Unidades del modelo: tonf, m, $^\circ$C. Eje X según los ejes 1--5, eje Y según los ejes A--B y eje Z vertical.
  \item Columnas, vigas y riostras se modelan con elementos tipo \emph{frame} ubicados en los ejes
        (no se consideran zonas rígidas ni excentricidades en las conexiones).
  \item Apoyos: columnas empotradas en la base (placas base con anclajes rígidos).
  \item Conexiones viga--columna rígidas (transmiten momento) en ambas direcciones.
  \item Riostras con extremos rotulados (liberación de $M_2$ y $M_3$ en ambos extremos), por lo que trabajan
        esencialmente a carga axial.
  \item Columnas orientadas con su eje fuerte resistiendo la flexión en el plano de los marcos rígidos
        (ejes 2, 3 y 4, dirección Y): ángulo local de $90^\circ$.
  \item La losa de H.A. de 18 cm se considera como diafragma rígido en su plano: se asigna una restricción
        tipo \emph{Diaphragm} (DIAF1 a DIAF5) a todos los nudos de cada nivel. La losa no se modela; sus cargas
        se aplican sobre las vigas (sección anterior).
  \item En cada nivel se crea un nudo en el centro de masas (nudos 199, 299, \dots, 599), incorporado al diafragma
        y restringido en $U_Z$, $R_X$ y $R_Y$, donde se aplican las fuerzas sísmicas.
  \item Patrones de carga: \texttt{PP} (peso propio $\times 1{,}10$), \texttt{SCP}, \texttt{L}, \texttt{LR},
        \texttt{EX} y \texttt{EY}. Combinaciones C1 a C11 y ENVOLVENTE.
  \item Análisis estático lineal de primer orden.
\end{itemize}

\subsection{Secciones de prediseño}
\begin{table}[H]
\centering
\caption{Secciones utilizadas en el modelo (propiedades calculadas por SAP2000 a partir de las dimensiones).}
\begin{tabular}{llcccc}
\toprule
Elemento & Perfil & $A$ [cm$^2$] & $I_{33}$ [cm$^4$] & $I_{22}$ [cm$^4$] & $Z_{33}$ [cm$^3$] \\
\midrule
Columnas (todas) & W14X90 & $168{,}5$ & $40917$ & $15019$ & $2527$ \\
Vigas ejes A y B & W21X50 & $93{,}4$ & $39988$ & $1037$ & $1766$ \\
Vigas ejes 1 a 5 & W21X62 & $116{,}4$ & $54560$ & $2391$ & $2333$ \\
Riostras ejes A y B & HSS8X8X1/2 & $90{,}4$ & $5541$ & $5541$ & $650$ \\
Riostras ejes 1 y 5 & HSS10X10X1/2 & $114{,}4$ & $11212$ & $11212$ & $1040$ \\
\bottomrule
\end{tabular}
\end{table}

Las secciones corresponden a un prediseño y serán verificadas/ajustadas en las Tareas 2 y 3.

\section{Resultados del análisis}

\subsection{Esfuerzos internos máximos}
La Tabla~\ref{tab:esf} resume los esfuerzos internos máximos (envolvente de las combinaciones LRFD) para cada
tipo de elemento, indicando el número de la barra en el modelo, su ubicación y la combinación que controla.
La numeración de barras se muestra en el Anexo~\ref{anx:num}. Por efecto del diafragma rígido, las vigas no
presentan esfuerzo axial en el modelo. Los diagramas de esfuerzos internos se adjuntan en el Anexo (capturas de SAP2000).

\begin{table}[H]
\centering
\small
\caption{Resumen de esfuerzos internos máximos de diseño (LRFD).}
\label{tab:esf}
\begin{tabular}{llcccc}
\toprule
Elemento & Esfuerzo & Valor & Barra & Ubicación & Comb. \\
\midrule
\multirow{5}{*}{Columnas} & $P_u$ compresión [tonf] & $280{,}3$ & 3 & Col. 3A, piso 1 & C2 \\
 & $P_u$ tracción [tonf] & $82{,}9$ & 1 & Col. 1A, piso 1 & C10 \\
 & $M_{u3}$ [tonf$\cdot$m] & $23{,}2$ & 18 & Col. 3B, piso 2 & C6 \\
 & $M_{u2}$ [tonf$\cdot$m] & $8{,}6$ & 12 & Col. 2A, piso 2 & C5 \\
 & $V_{u2}$ [tonf] & $15{,}0$ & 18 & Col. 3B, piso 2 & C6 \\
\midrule
\multirow{2}{*}{Vigas X (ejes A, B)} & $M_{u3}$ [tonf$\cdot$m] & $38{,}9$ & 52 & Eje A (2-3), nivel 1 & C2 \\
 & $V_{u2}$ [tonf] & $18{,}7$ & 52 & Eje A (2-3), nivel 1 & C2 \\
\midrule
\multirow{2}{*}{Vigas Y (ejes 1 a 5)} & $M_{u3}$ [tonf$\cdot$m] & $40{,}8$ & 93 & Eje 3 (A-B), nivel 1 & C7 \\
 & $V_{u2}$ [tonf] & $24{,}1$ & 93 & Eje 3 (A-B), nivel 1 & C2 \\
\midrule
\multirow{2}{*}{Riostras X (ejes A, B)} & $P_u$ compresión [tonf] & $79{,}5$ & 118 & Eje B (1-2), piso 1 & C5 \\
 & $P_u$ tracción [tonf] & $57{,}3$ & 119 & Eje B (4-5), piso 1 & C9 \\
\midrule
\multirow{2}{*}{Riostras Y (ejes 1, 5)} & $P_u$ compresión [tonf] & $108{,}0$ & 136 & Eje 1, piso 1 & C6 \\
 & $P_u$ tracción [tonf] & $102{,}4$ & 136 & Eje 1, piso 1 & C11 \\
\bottomrule
\end{tabular}
\end{table}

\subsection{Desplazamientos laterales}
A modo informativo, la Tabla~\ref{tab:desp} entrega los desplazamientos del centro de masas para las cargas
sísmicas \texttt{EX} y \texttt{EY} (sin mayorar) y las derivas de entrepiso, que resultan menores al límite
de $0{,}002$ de la NCh433.
\begin{table}[H]
\centering
\caption{Desplazamientos del CM y derivas de entrepiso (estados \texttt{EX} y \texttt{EY}).}
\label{tab:desp}
\begin{tabular}{ccccc}
\toprule
Nivel & $u_X$ [cm] & $\Delta_X/h$ [$\times 10^{-3}$] & $u_Y$ [cm] & $\Delta_Y/h$ [$\times 10^{-3}$] \\
\midrule
1 & $0{,}25$ & $0{,}747$ & $0{,}34$ & $1{,}021$ \\
2 & $0{,}49$ & $0{,}816$ & $0{,}71$ & $1{,}234$ \\
3 & $0{,}73$ & $0{,}789$ & $1{,}06$ & $1{,}173$ \\
4 & $0{,}94$ & $0{,}720$ & $1{,}38$ & $1{,}061$ \\
5 & $1{,}12$ & $0{,}592$ & $1{,}64$ & $0{,}860$ \\
\bottomrule
\end{tabular}
\end{table}

\section{Análisis plástico del eje 2}

\subsection{Modelo y supuestos}
Se estudia el marco rígido del eje 2 como un problema plano (plano YZ): un vano de $8{,}0$ m entre los ejes
A y B y cinco pisos, con columnas W14X90 (flexión en torno al eje fuerte) y vigas W21X62,
empotrado en la base. Se adoptan los siguientes supuestos:
\begin{itemize}
  \item Comportamiento elastoplástico perfecto, con rótulas plásticas concentradas en los extremos de vigas y
        columnas, de capacidad $M_p = Z_x F_y$. Se desprecia la interacción carga axial--momento y los efectos
        de segundo orden (P--$\Delta$), de forma de poder comparar directamente con el método cinemático.
  \item Cargas gravitacionales que tributan al eje 2 (aplicadas antes de la carga lateral):
        $D + 0{,}25L$, es decir, el peso propio de la viga ($\times 1{,}10$), la SCP y el $25\%$ de la
        sobrecarga de los paños 1--2 y 2--3 (carga trapezoidal más triangular sobre la viga), además de las reacciones de las
        vigas X que llegan a las columnas del eje.
  \item Carga lateral con la forma de la distribución sísmica de la Tabla~\ref{tab:sismo}: en el nivel $k$ se aplica
        $\alpha_k F$, con $\alpha_k = F_k/Q$, donde $F$ es la carga lateral total (corte basal del marco).
\end{itemize}

\begin{table}[H]
\centering
\begin{minipage}{0.45\textwidth}
\centering
\captionof{table}{Capacidades de las secciones.}
\begin{tabular}{lcc}
\toprule
 & Columna & Viga \\
 & W14X90 & W21X62 \\
\midrule
$S_x$ [cm$^3$] & $2301$ & $2046$ \\
$Z_x$ [cm$^3$] & $2527$ & $2333$ \\
$M_y = S_x F_y$ [tonf$\cdot$m] & $81{,}0$ & $72{,}0$ \\
$M_p = Z_x F_y$ [tonf$\cdot$m] & $88{,}9$ & $82{,}1$ \\
\bottomrule
\end{tabular}
\end{minipage}\hfill
\begin{minipage}{0.45\textwidth}
\centering
\captionof{table}{Patrón de carga lateral.}
\begin{tabular}{ccc}
\toprule
Nivel & $Z_k$ [m] & $\alpha_k = F_k/Q$ \\
\midrule
1 & $3{,}30$ & $0{,}1211$ \\
2 & $6{,}30$ & $0{,}1255$ \\
3 & $9{,}30$ & $0{,}1489$ \\
4 & $12{,}30$ & $0{,}1940$ \\
5 & $15{,}30$ & $0{,}4105$ \\
\bottomrule
\end{tabular}
\end{minipage}
\end{table}

\subsection{Curva de capacidad (análisis incremental, SAP2000)}
La curva de capacidad se obtiene mediante un análisis estático no lineal (\emph{pushover}) con control de
desplazamiento en el techo, partiendo del estado gravitacional. El procedimiento es evento a evento: se
incrementa la carga lateral hasta que alguna sección alcanza su $M_p$, se introduce una rótula en ese punto y se
continúa, hasta que la estructura se transforma en un mecanismo. La definición de las rótulas y de los casos
no lineales en SAP2000 se detalla en el Anexo~B.

\textbf{Primera fluencia.} Se define como el instante en que la fibra extrema de alguna sección alcanza $F_y$, es
decir, $|M_{grav} + F\,m_{lat}| = M_y$. Esto ocurre en la viga del nivel 2 para
\[
F_{y} = 47{,}2\tonf \qquad (\delta_{techo} = 10{,}17\ \mathrm{cm}).
\]

\textbf{Primera rótula plástica} ($M = M_p$): $F = 55{,}9\tonf$.

\textbf{Colapso.} El mecanismo se forma con la rótula N$^\circ$11 para
\[
F_{c} = 83{,}4\tonf \qquad (\delta_{techo} = 60{,}3\ \mathrm{cm}).
\]

\begin{figure}[H]
\centering
\begin{tikzpicture}
\begin{axis}[width=0.85\textwidth,height=7.5cm,grid=major,
  xlabel={Desplazamiento de techo $\delta$ [cm]},ylabel={Carga lateral total $F$ [tonf]},
  xmin=0,xmax=77,ymin=0,ymax=104,
  legend pos=south east,legend style={font=\small}]
\addplot[thick,blue,mark=*,mark size=1.5pt] coordinates {(0.013,0.000) (12.042,55.946) (12.348,57.047) (14.575,63.268) (17.014,68.601) (17.605,69.758) (21.883,74.415) (22.661,74.886) (26.896,76.359) (32.957,78.232) (59.771,83.369) (60.263,83.448) (75.329,83.448)};
\addlegendentry{Pushover (evento a evento)}
\addplot[red,dashed,thick] coordinates {(0,83.448) (76.84,83.448)};
\addlegendentry{Método cinemático $F_c = 83{,}4$ tonf}
\addplot[only marks,mark=square*,black,mark size=2.5pt] coordinates {(10.167,47.223)};
\addlegendentry{Primera fluencia $F_y = 47{,}2$ tonf}
\end{axis}
\end{tikzpicture}
\caption{Curva de capacidad del eje 2 (corte basal versus desplazamiento de techo).}
\label{fig:pushover}
\end{figure}

\begin{table}[H]
\centering
\caption{Secuencia de formación de rótulas plásticas.}
\begin{tabular}{clcc}
\toprule
N$^\circ$ & Ubicación & $F$ [tonf] & $\delta_{techo}$ [cm] \\
\midrule
1 & Viga nivel 2, extremo eje B & $55{,}9$ & $12{,}04$ \\
2 & Viga nivel 1, extremo eje B & $57{,}0$ & $12{,}35$ \\
3 & Viga nivel 3, extremo eje B & $63{,}3$ & $14{,}57$ \\
4 & Columna eje B, piso 1 (extremo inferior) & $68{,}6$ & $17{,}01$ \\
5 & Columna eje A, piso 1 (extremo inferior) & $69{,}8$ & $17{,}61$ \\
6 & Viga nivel 1, extremo eje A & $74{,}4$ & $21{,}88$ \\
7 & Viga nivel 2, extremo eje A & $74{,}9$ & $22{,}66$ \\
8 & Viga nivel 4, extremo eje B & $76{,}4$ & $26{,}90$ \\
9 & Viga nivel 3, extremo eje A & $78{,}2$ & $32{,}96$ \\
10 & Columna eje B, piso 3 (extremo superior) & $83{,}4$ & $59{,}77$ \\
11 & Columna eje A, piso 3 (extremo superior) & $83{,}4$ & $60{,}26$ \\
\bottomrule
\end{tabular}
\end{table}
En la tabla, ``extremo eje A/B'' indica el extremo de la viga donde se forma la rótula. Algunas de las rótulas
formadas en la secuencia se descargan una vez que se forma el mecanismo definitivo.

\subsection{Método cinemático}
El método cinemático (teorema del límite superior) iguala el trabajo externo de las cargas con el trabajo interno
disipado en las rótulas para un mecanismo cinemáticamente admisible con giro virtual $\theta$:
\[
W_e = \sum_k \alpha_k F\,\delta_k = W_i = \sum_r M_{p,r}\,\theta_r
\quad\Rightarrow\quad
F = \frac{\sum_r M_{p,r}\,\theta_r}{\sum_k \alpha_k\,\delta_k}.
\]
Se evalúa la familia de mecanismos de piso y combinados en que los pisos $j$ a $m$ se desplazan lateralmente
(columnas de esos pisos giran $\theta$), con rótulas: en el pie de las columnas del piso $j$, en ambos extremos de
las vigas de los niveles $j$ a $m-1$ (o en las columnas si son más débiles, $\min(M_{pb},\,2M_{pc})$ por nudo),
y en la cabeza de las columnas del piso $m$ (o en las vigas de techo si $m = 5$). Para $j = m$ se obtiene el
mecanismo de piso blando, y para $j=1,\ m=5$ el mecanismo global de vigas. Como el mecanismo solo
involucra traslaciones horizontales de las vigas, las cargas gravitacionales no realizan trabajo.

Por ejemplo, para el mecanismo que controla (pisos 1 a 3):
\[
W_i = \left[4\,M_{pc} + 4\,M_{pb}\right]\theta
     = 684{,}0\,\theta \ \mathrm{tonf\cdot m},
\qquad
W_e = F\,\theta\sum_k \alpha_k\,\delta_k/\theta = 8{,}197\,F\,\theta\ \mathrm{m}
\]
\[
\Rightarrow\quad F_c = \frac{684{,}0}{8{,}197} = 83{,}4\tonf .
\]

\begin{table}[H]
\centering
\caption{Cargas de colapso de los mecanismos evaluados (los 8 menores).}
\begin{tabular}{lccc}
\toprule
Mecanismo & $W_i/\theta$ [tonf$\cdot$m] & $W_e/(F\theta)$ [m] & $F$ [tonf] \\
\midrule
Combinado pisos 1 a 3 & $684{,}0$ & $8{,}197$ & $83{,}4$ \\
Combinado pisos 1 a 4 & $848{,}2$ & $10{,}010$ & $84{,}7$ \\
Combinado pisos 1 a 2 & $519{,}8$ & $5{,}937$ & $87{,}6$ \\
Global (vigas + base) & $998{,}7$ & $11{,}242$ & $88{,}8$ \\
Combinado pisos 2 a 4 & $684{,}0$ & $6{,}710$ & $101{,}9$ \\
Combinado pisos 2 a 5 & $834{,}5$ & $7{,}942$ & $105{,}1$ \\
Combinado pisos 2 a 3 & $519{,}8$ & $4{,}897$ & $106{,}2$ \\
Piso 1 (piso blando) & $355{,}7$ & $3{,}300$ & $107{,}8$ \\
\bottomrule
\end{tabular}
\end{table}

\begin{figure}[H]
\centering
\begin{tikzpicture}[x=0.33cm,y=0.33cm]
\draw[thick] (0,0) -- (0,15.3);
\draw[pattern=north east lines] (0,0) ++(-0.8,-0.5) rectangle ++(1.6,0.5);
\draw[thick] (8,0) -- (8,15.3);
\draw[pattern=north east lines] (8,0) ++(-0.8,-0.5) rectangle ++(1.6,0.5);
\draw[thick] (0,3.3) -- (8,3.3);
\draw[thick] (0,6.3) -- (8,6.3);
\draw[thick] (0,9.3) -- (8,9.3);
\draw[thick] (0,12.3) -- (8,12.3);
\draw[thick] (0,15.3) -- (8,15.3);
\draw[-{Stealth},blue,thick] (-2.22681,3.3) -- (0,3.3) node[midway,above,font=\tiny]{$0{,}121 F$};
\draw[-{Stealth},blue,thick] (-2.25301,6.3) -- (0,6.3) node[midway,above,font=\tiny]{$0{,}126 F$};
\draw[-{Stealth},blue,thick] (-2.39321,9.3) -- (0,9.3) node[midway,above,font=\tiny]{$0{,}149 F$};
\draw[-{Stealth},blue,thick] (-2.66406,12.3) -- (0,12.3) node[midway,above,font=\tiny]{$0{,}194 F$};
\draw[-{Stealth},blue,thick] (-3.96291,15.3) -- (0,15.3) node[midway,above,font=\tiny]{$0{,}410 F$};
\fill[red] (0,0.5) circle (0.35); \fill[red] (8,0.5) circle (0.35);
\fill[red] (0.6,3.3) circle (0.35); \fill[red] (7.4,3.3) circle (0.35);
\fill[red] (0.6,6.3) circle (0.35); \fill[red] (7.4,6.3) circle (0.35);
\fill[red] (0,8.8) circle (0.35); \fill[red] (8,8.8) circle (0.35);
\node[below] at (0,-0.8) {A}; \node[below] at (8,-0.8) {B};
\end{tikzpicture}
\caption{Mecanismo de colapso del eje 2 (rótulas plásticas en rojo).}
\label{fig:mecanismo}
\end{figure}

\subsection{Comparación y discusión}
\begin{table}[H]
\centering
\caption{Comparación de resultados del análisis plástico del eje 2.}
\begin{tabular}{lc}
\toprule
Parámetro & Valor \\
\midrule
Carga de primera fluencia $F_y$ & $47{,}2\tonf$ \\
Carga de primera rótula plástica & $55{,}9\tonf$ \\
Carga de colapso, pushover $F_c$ & $83{,}4\tonf$ \\
Carga de colapso, método cinemático & $83{,}4\tonf$ \\
Diferencia & $0{,}00\%$ \\
Sobrerresistencia $F_c/F_y$ & $1{,}77$ \\
\bottomrule
\end{tabular}
\end{table}

Ambos métodos entregan la misma carga de colapso, ya que el método cinemático es un límite superior que
coincide con la solución exacta cuando se evalúa el mecanismo correcto, y el análisis incremental con
rótulas elastoplásticas perfectas converge al mismo mecanismo. Las diferencias que se observen respecto al
resultado de SAP2000 se deben a: (i) la definición de las rótulas (si se utilizan rótulas automáticas
ASCE~41 con endurecimiento, la curva no presenta una meseta perfectamente horizontal), (ii) la
interacción P--M en las columnas si se usan rótulas tipo P--M2--M3 y (iii) los efectos P--$\Delta$,
que reducen la capacidad lateral. El momento máximo de tramo de las vigas bajo cargas gravitacionales es
$17{,}4\tonfm < M_{pb}$, por lo que no se forman rótulas en el interior de las vigas
y es válido considerar rótulas solo en los extremos.
La razón $F_c/F_y = 1{,}77$ refleja la reserva de resistencia del
marco hiperestático por redistribución de momentos luego de la primera fluencia.

\section{Diseño de la riostra más traccionada}

\subsection{Demanda}
Del análisis, la riostra más traccionada corresponde a la barra 136 del modelo,
ubicada en el \textbf{eje 1, piso 1}, que une el nudo 101
$(X,Y,Z) = (0{,}0;0{,}0;3{,}3)$ con el nudo 11
$(0{,}0;8{,}0;0{,}0)$ (ver Figuras~\ref{fig:planta} y \ref{fig:elev}).
Su largo es $L = 8{,}65$ m y la combinación que controla es C11 ($0.9D - 1.4E_y$):
\[
T_u = 102{,}4\tonf .
\]
La riostra simétrica del eje 5 (mismo piso) presenta una demanda equivalente.

\subsection{Sección propuesta}
Se propone un perfil tubular cuadrado \textbf{HSS10X10X1/2} de acero ASTM A500 Gr.C
($F_y = 3518\ \mathrm{kgf/cm^2}$, $F_u = 4354\ \mathrm{kgf/cm^2}$), con espesor de diseño $t = 0{,}93\,t_{nom} = 1{,}18$ cm, $A_g = 111{,}2\ \mathrm{cm^2}$ y
$r = 9{,}80$ cm.

\subsection{Verificaciones (AISC 360-16, Capítulo D)}

\textbf{Esbeltez (recomendación D1).}
\[
\frac{L}{r} = \frac{865{,}4}{9{,}80} = 88{,}3 \le 300 \quad \checkmark
\]

\textbf{Fluencia en el área bruta (D2-a), $\phi_t = 0{,}90$.}
\[
\phi_t P_n = 0{,}90 \cdot F_y \cdot A_g = 0{,}90 \cdot 3518 \cdot 111{,}2
= 352{,}0\tonf
\]

\textbf{Ruptura en el área neta efectiva (D2-b), $\phi_t = 0{,}75$.}
La conexión se materializa con una plancha gusset concéntrica de espesor $t_g = 2{,}2$ cm insertada en
una ranura del tubo de ancho $t_g + 2$ mm, y soldada con cuatro cordones de filete de largo $l$.
\[
A_n = A_g - 2\,t\,(t_g + 0{,}2) = 111{,}2 - 2\cdot1{,}18\cdot(2{,}2+0{,}2)
= 105{,}5\ \mathrm{cm^2}
\]
Factor de corte diferido (Tabla D3.1, caso 6, HSS rectangular con una plancha concéntrica, $l \ge H$):
\[
\bar{x} = \frac{B^2 + 2BH}{4(B+H)} = 9{,}53\ \mathrm{cm},
\qquad
U = 1 - \frac{\bar{x}}{l} = 1 - \frac{9{,}53}{30} = 0{,}682
\]
\[
A_e = U A_n = 72{,}0\ \mathrm{cm^2}
\qquad\Rightarrow\qquad
\phi_t P_n = 0{,}75 \cdot 4354 \cdot 72{,}0 = 235{,}2\tonf
\]

\textbf{Resistencia de diseño.}
\[
\phi_t P_n = \min(352{,}0;\ 235{,}2) = 235{,}2\tonf
\ \ge\ T_u = 102{,}4\tonf \quad \checkmark
\qquad \text{F.U.} = 0{,}44
\]

\subsection{Conexión preliminar}
Se propone una conexión soldada con plancha gusset de acero ASTM A36, $t_g = 22$ mm, insertada
en el tubo ranurado y unida mediante cuatro cordones de soldadura de filete E70XX.

\textbf{Soldadura.} Tamaño de filete $w = 8$ mm (mínimo según Tabla J2.4: 5 mm para
$6 < t \le 13$ mm; máximo $t - 2 = 9{,}8$ mm).
Resistencia por unidad de largo de un cordón (J2.4), $\phi = 0{,}75$:
\[
\phi R_n/l = 0{,}75 \cdot 0{,}60 F_{EXX} \cdot 0{,}707\,w
= 0{,}75 \cdot 0{,}60 \cdot 4920 \cdot 0{,}707 \cdot 0{,}8
= 1252\ \mathrm{kgf/cm}
\]
Corte del metal base (pared del tubo): $0{,}75\cdot 0{,}60\,F_u\,t = 2314\ \mathrm{kgf/cm}$ (no controla).
Largo requerido de cada cordón:
\[
l_{req} = \frac{T_u}{4\,\phi R_n/l} = \frac{102378}{4 \cdot 1252}
= 20{,}4\ \mathrm{cm}
\quad\Rightarrow\quad l = 30\ \mathrm{cm} \ (\ge B = 25{,}4\ \mathrm{cm})
\]

\textbf{Plancha gusset.} Ancho de Whitmore $L_w = B + 2\,l\tan 30^\circ = 60{,}0$ cm;
área de corte $A_{gv} = 2\,l\,t_g = 132{,}0\ \mathrm{cm^2}$; área en tracción $A_{nt} = B\,t_g = 55{,}9\ \mathrm{cm^2}$.

\begin{table}[H]
\centering
\caption{Resumen de verificaciones de la riostra traccionada y su conexión.}
\begin{tabular}{llcc}
\toprule
Estado límite & Expresión & $\phi R_n$ [tonf] & $T_u/\phi R_n$ \\
\midrule
Fluencia en área bruta & $\phi_t F_y A_g$ & $352{,}0$ & $0{,}29$ \\
Ruptura en área neta efectiva & $\phi_t F_u A_e$ & $235{,}2$ & $0{,}44$ \\
Soldaduras (4 cordones) & $4\,l\,\phi R_n/l$ & $150{,}3$ & $0{,}68$ \\
Plancha gusset: fluencia (Whitmore) & $\phi F_y L_w t_g$ & $300{,}6$ & $0{,}34$ \\
Plancha gusset: bloque de corte & $\phi R_{bs}$ & $321{,}2$ & $0{,}32$ \\
Plancha gusset: ruptura por corte & $\phi\,0{,}6F_u\,2l\,t_g$ & $242{,}3$ & $0{,}42$ \\
\bottomrule
\end{tabular}
\end{table}

\begin{figure}[H]
\centering
\begin{tikzpicture}[x=0.095cm,y=0.095cm,font=\small]
% gusset
\draw[thick,fill=gray!15] (-42,-30) -- (40,-30) -- (40,30) -- (-42,30) -- cycle;
\node[align=center] at (-5,-38) {Gusset $t_g = 22$ mm (ASTM A36)};
% tubo
\draw[thick,fill=white] (0,12.7) -- (90,12.7) -- (90,-12.7) -- (0,-12.7) -- cycle;
\draw[dashed] (0,0) -- (90,0);
\node at (28,0) [above] {HSS10X10X1/2};
% soldaduras
\draw[red,line width=2pt] (0,12.7) -- (30,12.7);
\draw[red,line width=2pt] (0,-12.7) -- (30,-12.7);
\draw[<->] (0,17.7) -- node[above]{$l = 30$ cm} (30,17.7);
\draw[<->] (82,-12.7) -- node[left,fill=white,inner sep=1pt]{$B = 25{,}4$ cm} (82,12.7);
\node[red] at (15,-17.7) {4 filetes $w = 8$ mm, E70XX};
% Whitmore
\draw[blue,dashed] (0,12.7) -- (-30,30.0205);
\draw[blue,dashed] (0,-12.7) -- (-30,-30.0205);
\draw[blue] (-30,-30.0205) -- node[left]{$L_w$} (-30,30.0205);
\draw[->,very thick] (92,0) -- (108,0) node[right]{$T_u$};
\end{tikzpicture}
\caption{Esquema de la conexión preliminar riostra--gusset (vista en elevación de la riostra).}
\label{fig:conexion}
\end{figure}

\textbf{Comentario.} La sección queda con un factor de utilización bajo a tracción; sin embargo, la misma riostra
trabaja en compresión al invertirse el sentido del sismo ($P_u \approx T_u$), condición que controlará su diseño
en la Tarea 2, por lo que se mantiene el perfil HSS10X10X1/2. La razón ancho--espesor del tubo es
$b/t = 18{,}5$.

\section{Conclusiones}
\begin{itemize}
  \item El peso sísmico de la estructura es $P = 1014{,}2\tonf$, lo que resulta en un corte
        basal $Q = 182{,}6\tonf$, aplicado en el centro de masas de cada piso con la distribución
        en altura indicada.
  \item Los marcos arriostrados concentran la mayor parte de la carga lateral; las riostras de los ejes 1 y 5
        (dirección Y) son las más solicitadas, con $T_u = 102{,}4\tonf$.
  \item El análisis plástico del eje 2 entrega una carga lateral de primera fluencia
        $F_y = 47{,}2\tonf$, una primera rótula plástica para $F = 55{,}9\tonf$ y
        una carga de colapso $F_c = 83{,}4\tonf$. El método cinemático entrega
        $F_c = 83{,}4\tonf$ (mecanismo \emph{Combinado pisos 1 a 3}), coincidente con el
        análisis incremental, ya que este último corresponde a la solución exacta del problema rígido-plástico
        con rótulas en los extremos de los elementos.
  \item La riostra más traccionada se diseña con un perfil HSS10X10X1/2 (ASTM A500 Gr.C) con un factor de utilización
        de $0{,}44$ a tracción. La sección queda holgada en tracción, pero se escogió anticipando el
        diseño a compresión de la Tarea 2 (la misma riostra se comprime al invertir el sismo).
\end{itemize}

\appendix
\section{Numeración del modelo}\label{anx:num}
\begin{itemize}
  \item \textbf{Nudos:} $100k + 10\,i_Y + i_X$, con $k$ = nivel (0 = base), $i_Y = 0$ (eje A) o 1 (eje B) e
        $i_X = 1 \dots 5$ (ejes 1 a 5). Ejemplo: nudo 312 = nivel 3, eje B, eje 2. Los nudos de centro de masas
        son $100k + 99$.
  \item \textbf{Barras:} 1--50 columnas (10 por piso, piso 1 = barras 1--10), 51--90 vigas X (ejes A y B),
        91--115 vigas Y (ejes 1 a 5), 116--135 riostras X (ejes A y B) y 136--145 riostras Y (ejes 1 y 5).
  \item \textbf{Grupos} definidos en SAP2000 (creados por el script de la API): COLUMNAS, VIGAS\_X, VIGAS\_Y, RIOSTRAS\_X, RIOSTRAS\_Y y EJE\_2.
\end{itemize}

\section{Análisis no lineal en SAP2000 (eje 2)}
Se utiliza el modelo 2D \texttt{G5\_Eje2\_Pushover} (SAP2000 v27, plano YZ), que contiene la geometría, las secciones,
las cargas gravitacionales tributarias del eje 2 (patrones \texttt{PP}, \texttt{SCP}, \texttt{L}, \texttt{LR})
y el patrón lateral \texttt{PUSH} (fuerzas $\alpha_k$ en cada nivel, de suma unitaria). Los pasos son:
\begin{enumerate}
  \item \emph{Define $\rightarrow$ Section Properties $\rightarrow$ Hinge Properties}: crear rótulas
        \texttt{RC} (columna) y \texttt{RV} (viga) tipo \emph{Moment M3}, comportamiento elastoplástico perfecto
        (\emph{Deformation Controlled}), con $M_p = 88{,}9\tonfm$ y
        $M_p = 82{,}1\tonfm$ respectivamente, y meseta horizontal (puntos B--C con igual momento).
        Alternativamente, rótulas automáticas \emph{Auto M3} según ASCE 41-13 Tabla 9-6.
  \item Seleccionar todas las barras y asignar (\emph{Assign $\rightarrow$ Frame $\rightarrow$ Hinges}) rótulas a
        distancias relativas $0$ y $1$.
  \item \emph{Define $\rightarrow$ Load Cases}: caso \texttt{CGNL}, \emph{Static Nonlinear}, cargas
        $PP + SCP + 0{,}25L + 0{,}25L_r$, control de carga, condición inicial cero.
  \item Caso \texttt{PUSHOVER}, \emph{Static Nonlinear}, continuar desde \texttt{CGNL}, carga \texttt{PUSH} con
        factor 1, control de desplazamiento (\emph{Monitored displacement}: nudo 51, dirección U2, $60$--$80$ cm),
        resultados en múltiples pasos.
  \item Ejecutar y obtener la curva en \emph{Display $\rightarrow$ Show Static Pushover Curve} (corte basal vs.
        desplazamiento del nudo de techo).
\end{enumerate}

\section*{Anexo C: Diagramas de esfuerzos internos}
\addcontentsline{toc}{section}{Anexo C: Diagramas de esfuerzos internos}
% Insertar aquí capturas de SAP2000, por ejemplo:
% \begin{figure}[H]\centering
%   \includegraphics[width=0.9\textwidth]{diagrama_M3_eje2.png}
%   \caption{Diagrama de momento $M_3$, eje 2, combinación ENVOLVENTE.}
% \end{figure}
Ver capturas del modelo SAP2000 (diagramas de momento, corte y axial para la combinación ENVOLVENTE).

\end{document}
